Analizuję problem Online Bipartite Matching - problem dobrze znany i przebadany z powodu jego użyteczności. Jest on pomocny w modelowaniu różnych realistycznych problemów w modelach powiązanych ze sobą. Moim celem jest stwierdzenie jakie modyfikacje modelu pozwalają deterministycznym algorytmom osiągnąć miarę competitiveness ponad $0.5$. Ten problem wprowadzili Richard M. Karp, Umesh V. Vazirani i Vijay V. Vazirani. Analiza została przeprowadzona przy różnych założeniach: ograniczenie adwersarza przez wymuszenie losowej kolejności ujawniania wierzchołków lub wzmacnianie algorytmu poprzez przekazywanie informacji o ostatecznym stopniu wierzchołków albo konstrukcji skojarzenia ułamkowego.